\documentclass[11pt]{article}
\usepackage[margin=1in]{geometry}
\usepackage{amsmath,amssymb,amsthm,mathtools}
\usepackage{graphicx}
\usepackage{microtype}
\usepackage{titlesec}
\usepackage[hidelinks]{hyperref}

\titleformat{\section}{\normalfont\Large\bfseries}{\thesection}{1em}{}

\newtheorem{theorem}{Theorem}
\newcommand{\R}{\mathbb{R}}
\newcommand{\Sph}{\mathbb{S}}
\newcommand{\ip}[2]{\langle #1,#2\rangle}

\title{Legendrian Necks of the Positive Pseudosphere\\
Are Convex Normal Lifts}
\author{Candidate full solution packet for arXiv:math/0610904}
\date{21 July 2026}

\begin{document}
\maketitle

\begin{center}
\fbox{\parbox{0.9\textwidth}{
\textbf{Status.} Candidate full solution, likely valid, pending human review.
The theorem below answers the Legendrian-neck conjecture in Section~5.1 of
Kuperberg's paper for the smooth spacelike notion of neck used there.
}}
\end{center}

\section{Source problem}

Let
\[
 H^+=\{(x,y)\in\R^n\times\R^n:\ip{x}{y}=1\}.
\]
Kuperberg equips the ambient space with the split bilinear form
\[
 \big\langle (\xi,\eta),(\xi',\eta')\big\rangle_Q
 =\frac12\big(\ip{\xi}{\eta'}+\ip{\xi'}{\eta}\big).
\]
If a convex body $K\subset\R^n$ contains the origin in its interior, its
positive normal lift is
\[
 K^+=\{(x,y)\in\partial K\times\partial K^\circ:\ip{x}{y}=1\}.
\]
On source PDF page~9, Section~5.1, the paper conjectures that every
Legendrian neck of $H^+$ is $K^+$ for a convex body $K$ \cite{Kuperberg}.
The complete page is reproduced in the source-page evidence section.

The standard contact form on $H^+$, up to a nonzero scalar and sign, is
\[
 \alpha=\frac12\big(\ip{y}{dx}-\ip{x}{dy}\big).
\]
A neck is, in particular, a compact connected spacelike submanifold of
dimension $n-1$ which is a section in Kuperberg's topological sense. That
section condition gives degree one for its projection to the sphere of
$y$-directions.

\section{Proof intuition}

A point $(x,y)$ of a convex normal lift records a boundary point $x$ and its
supporting covector $y$. The conjecture asks whether the contact condition and
spacelikeness already force this convex interpretation.

They do. Legendrianity says that infinitesimal motion in $x$ is orthogonal to
$y$, so $y$ behaves like a normal. The neck condition makes the normal
direction a global parameter $u\in\Sph^{n-1}$. Writing $y=u/h(u)$, the contact
equations recover the usual inverse Gauss parametrization
$x=\nabla_{\Sph}h+hu$. The split metric on the neck becomes exactly the
spherical support tensor $h^{-1}(\nabla_{\Sph}^2h+hI)$. Its positivity is the
strict-convexity condition. A geodesic ODE proves the global support inequality
directly, so no local-to-global convexity theorem is needed.

\section{\hspace{0.75em}The theorem}

\begin{theorem}
Let $n\ge1$, and let $N\subset H^+$ be a smooth spacelike neck in the sense of
Kuperberg. If $N$ is Legendrian for $\ker\alpha$, then there is a unique smooth
strictly convex body $K\subset\R^n$, with $0\in\operatorname{int}K$, such that
$N=K^+$.
\end{theorem}

\begin{proof}
For $n=1$, a neck consists of one point $(a,a^{-1})$ on each of the two
components of $H^+$. If the first coordinates are $a_-<0<a_+$, it is exactly
$[a_-,a_+]^+$, so the assertion is immediate. Assume henceforth that $n\ge2$.

Write a tangent vector to $N$ at $(x,y)$ as $(\dot x,\dot y)$. Since $N$ lies
in the hypersurface $\ip{x}{y}=1$,
\begin{equation}\label{eq:level}
 \ip{y}{\dot x}+\ip{x}{\dot y}=0.
\end{equation}
Since $N$ is Legendrian,
\begin{equation}\label{eq:legendrian}
 \ip{y}{\dot x}-\ip{x}{\dot y}=0.
\end{equation}
Adding and subtracting \eqref{eq:level} and \eqref{eq:legendrian} gives
\begin{equation}\label{eq:orthogonality}
 \ip{y}{\dot x}=0,
 \qquad
 \ip{x}{\dot y}=0.
\end{equation}

Consider the normal-direction map
\[
 \nu:N\longrightarrow\Sph^{n-1},
 \qquad \nu(x,y)=\frac{y}{|y|}.
\]
It is a local diffeomorphism. Indeed, if $d\nu(\dot x,\dot y)=0$, then
$\dot y=a y$ for some scalar $a$. The second identity in
\eqref{eq:orthogonality} gives $a=\ip{x}{\dot y}=0$. Thus $\dot y=0$, and the
split squared norm of the tangent vector is
\[
 \big\langle(\dot x,\dot y),(\dot x,\dot y)\big\rangle_Q
 =\ip{\dot x}{\dot y}=0.
\]
Spacelikeness then forces $(\dot x,\dot y)=0$.

For completeness, the ambient projection
\[
 H^+\longrightarrow\Sph^{n-1},\qquad (x,y)\longmapsto y/|y|
\]
has contractible fibers: over $u$ one may write
$y=r u$ and $x=r^{-1}u+z$, with $r>0$ and $z\perp u$. Hence it is a homotopy
equivalence. The section condition in the definition of a neck implies that
$\nu$ has degree of absolute value one. A local diffeomorphism from the compact
connected neck to the sphere is a finite covering; its number of sheets is the
absolute degree. Consequently $\nu$ is a diffeomorphism. This also covers the
$n=2$ case, where multiple coverings of the circle must explicitly be ruled
out.

We may therefore parameterize $N$ by $u\in\Sph^{n-1}$ and write
\[
 y(u)=\frac{u}{h(u)},
 \qquad h(u)>0.
\]
The equation $\ip{x}{y}=1$ says $\ip{x(u)}{u}=h(u)$. Moreover,
\eqref{eq:orthogonality} gives $\ip{u}{dx}=0$. Differentiating
$h=\ip{x}{u}$ in a tangent direction $w\in T_u\Sph^{n-1}$ therefore yields
\[
 dh(w)=\ip{x}{w}.
\]
The tangential component of $x$ is thus $\nabla_{\Sph}h$, while its normal
component is $hu$. Hence
\begin{equation}\label{eq:xparam}
 x(u)=\nabla_{\Sph}h(u)+h(u)u.
\end{equation}

Put
\[
 A_u=\nabla_{\Sph}^2h(u)+h(u)I
 \quad\text{on }T_u\Sph^{n-1}.
\]
Using the Euclidean derivative of the spherical gradient in
\eqref{eq:xparam}, the two normal terms cancel and give
\[
 dx(w)=A_uw.
\]
Also
\[
 dy(w)=h^{-1}w-h^{-2}dh(w)u.
\]
Since $dx(w)$ is tangent to the sphere and $A_u$ is self-adjoint, the metric
induced on $N$ is
\begin{equation}\label{eq:metric}
 g_N(w,z)
 =\frac12\big(\ip{dx(w)}{dy(z)}+\ip{dx(z)}{dy(w)}\big)
 =h(u)^{-1}\ip{A_uw}{z}.
\end{equation}
Spacelikeness and $h>0$ imply that $A_u$ is positive definite for every $u$.

We now prove directly that $h$ is a global support function. Define
\[
 K=\bigcap_{v\in\Sph^{n-1}}
 \{z\in\R^n:\ip{z}{v}\le h(v)\}.
\]
Because $h$ is positive and continuous, $K$ is a compact convex body with the
origin in its interior. Fix $u,v\in\Sph^{n-1}$ with $v\ne u$, and let
$\gamma:[0,L]\to\Sph^{n-1}$ be a shortest unit-speed geodesic from $u$ to
$v$, where $0<L\le\pi$. Set
\[
 F(t)=h(\gamma(t))-\ip{x(u)}{\gamma(t)}.
\]
Equation \eqref{eq:xparam} gives $F(0)=F'(0)=0$. Since a Euclidean linear
function restricted to a great circle satisfies $\ell''+\ell=0$, we have
\[
 F''(t)+F(t)
 =\ip{A_{\gamma(t)}\dot\gamma(t)}{\dot\gamma(t)}>0.
\]
Variation of constants gives
\[
 F(t)=\int_0^t\sin(t-s)
 \ip{A_{\gamma(s)}\dot\gamma(s)}{\dot\gamma(s)}\,ds.
\]
For $0<t\le\pi$ this is strictly positive. In particular,
\[
 \ip{x(u)}{v}<h(v)\quad(v\ne u),
 \qquad
 \ip{x(u)}{u}=h(u).
\]
Thus $x(u)\in K$ and is the unique support point of $K$ with outward normal
$u$. It follows that $h$ is the support function of $K$ and
$\partial K=\{x(u):u\in\Sph^{n-1}\}$. The positive definiteness of $A_u$
shows that this boundary is smooth and strictly convex.

Finally, $y(u)=u/h(u)$ lies in $K^\circ$, because
$\ip{z}{u}\le h(u)$ for every $z\in K$, and equality holds at $x(u)$. Hence
$y(u)\in\partial K^\circ$ and
$\ip{x(u)}{y(u)}=1$. We have therefore shown
\[
 N=\{(x(u),y(u)):u\in\Sph^{n-1}\}=K^+.
\]
The support function $h(u)=1/|y(u)|$ is determined by $N$, so $K$ is unique.
\end{proof}

\section{\hspace{0.75em}Verification notes}

\begin{itemize}
\item The sign of the contact form is immaterial. The level-set equation and
Legendrian equation independently give the two vanishing pairings in
\eqref{eq:orthogonality}.
\item The only global topological input is already part of the source's notion
of a neck: its sphere-direction projection has degree one. This is needed in
dimension two to exclude a multiple covering.
\item Formula \eqref{eq:metric} checks the decisive sign: spacelike means
$\nabla_{\Sph}^2h+hI$ is positive, exactly the strict support-function tensor.
\item The global support inequality is proved by the displayed geodesic ODE;
no unquoted Hadamard or support-function realization theorem is used.
\item The proof concerns smooth spacelike necks, exactly the class in the
source conjecture. It does not claim the analogous statement for weak necks
with null tangents.
\end{itemize}

\section{\hspace{0.75em}Novelty check and limitations}

On 21 July 2026, a bounded search covered the run registry, solution, attempt,
and proof-gap indexes; the local source corpus; and web searches using the
exact phrases ``every Legendrian neck'', ``Legendrian neck convex body'', and
the source title together with ``Legendrian''. The search returned the source
paper and mirrors but no later paper answering this exact conjecture. This
supports, but does not prove, novelty. The result should remain a candidate
full solution until a specialist checks the contact convention and the
degree-one consequence of the neck definition.

\clearpage
\section{Source-page evidence}
\begin{center}
\centering
\includegraphics[width=\textwidth,height=0.88\textheight,keepaspectratio]
{figures/open_problem_crop.png}

\smallskip
\small Source PDF page 9. The target conjecture begins at the bottom of the
left column and concludes at the top of the right column.
\end{center}

\begin{thebibliography}{9}
\bibitem{Kuperberg}
G. Kuperberg,
\emph{From the Mahler conjecture to Gauss linking integrals},
Geom. Funct. Anal. \textbf{18} (2008), 870--892;
\href{https://arxiv.org/abs/math/0610904}{arXiv:math/0610904},
DOI: \href{https://doi.org/10.1007/s00039-008-0669-4}
{10.1007/s00039-008-0669-4}.
\end{thebibliography}

\end{document}
